%%% Hlavní soubor. Zde se definují základní parametry a odkazuje se na ostatní části. %%%

%% Verze pro jednostranný tisk:
% Okraje: levý 40mm, pravý 25mm, horní a dolní 25mm
% (ale pozor, LaTeX si sám přidává 1in)
\documentclass[12pt,a4paper]{report}
\setlength\textwidth{145mm}
\setlength\textheight{247mm}
\setlength\oddsidemargin{15mm}
\setlength\evensidemargin{15mm}
\setlength\topmargin{0mm}
\setlength\headsep{0mm}
\setlength\headheight{0mm}
% \openright zařídí, aby následující text začínal na pravé straně knihy
\let\openright=\clearpage

%% Pokud tiskneme oboustranně:
% \documentclass[12pt,a4paper,twoside,openright]{report}
% \setlength\textwidth{145mm}
% \setlength\textheight{247mm}
% \setlength\oddsidemargin{15mm}
% \setlength\evensidemargin{0mm}
% \setlength\topmargin{0mm}
% \setlength\headsep{0mm}
% \setlength\headheight{0mm}
% \let\openright=\cleardoublepage

\def\nazevprace{Název práce} % TODO

%% csLaTeX
\usepackage{czech}

\usepackage[utf8]{inputenc}

\usepackage{graphicx}
\usepackage{amsthm}

%% Balíček hyperref, kterým jdou vyrábět klikací odkazy v PDF,
%% ale hlavně ho používáme k uložení metadat do PDF (včetně obsahu).
%% POZOR, nezapomeňte vyplnit jméno práce a autora.
\usepackage[ps2pdf,unicode]{hyperref}   % Musí být za všemi ostatními balíčky
\hypersetup{pdftitle=Název práce} % TODO
\hypersetup{pdfauthor=Michael Pokorný}

%%% Drobné úpravy stylu

% Tato makra přesvědčují mírně ošklivým trikem LaTeX, aby hlavičky kapitol
% sázel příčetněji a nevynechával nad nimi spoustu místa. Směle ignorujte.
\makeatletter
\def\@makechapterhead#1{
  {\parindent \z@ \raggedright \normalfont
   \Huge\bfseries \thechapter. #1
   \par\nobreak
   \vskip 20\p@
}}
\def\@makeschapterhead#1{
  {\parindent \z@ \raggedright \normalfont
   \Huge\bfseries #1
   \par\nobreak
   \vskip 20\p@
}}
\makeatother

% Toto makro definuje kapitolu, která není očíslovaná, ale je uvedena v obsahu.
\def\chapwithtoc#1{
\chapter*{#1}
\addcontentsline{toc}{chapter}{#1}
}

\begin{document}

% Trochu volnější nastavení dělení slov, než je default.
\lefthyphenmin=2
\righthyphenmin=2

%%% Titulní strana práce
\pagestyle{empty}
\begin{center}

\large

Univerzita Karlova v Praze

\medskip

Matematicko-fyzikální fakulta

\vfill

{\bf\Large BAKALÁŘSKÁ PRÁCE}

\vfill

\centerline{\mbox{\includegraphics[width=60mm]{img/logo.eps}}}

\vfill
\vspace{5mm}

{\LARGE Michael Pokorný}

\vspace{15mm}

{\LARGE\bfseries Praktické datové struktury}

\vfill

Katedra aplikované matematiky

\vfill

\begin{tabular}{rl}

Vedoucí bakalářské práce: & Mgr. Martin Mareš, Ph.D.\\
\noalign{\vspace{2mm}}
Studijní program: & informatika \\
\noalign{\vspace{2mm}}
Studijní obor: & obecná informatika \\
\end{tabular}

\vfill

Praha 2014

\end{center}

\newpage

%%% Následuje vevázaný list -- kopie podepsaného "Zadání diplomové práce".
%%% Toto zadání NENÍ součástí elektronické verze práce, nescanovat.

%%% Na tomto místě mohou být napsána případná poděkování (vedoucímu práce,
%%% konzultantovi, tomu, kdo zapůjčil software, literaturu apod.)

\openright

\noindent
TODO: Poděkování.

\newpage
%%% Strana s čestným prohlášením k bakalářské práci

\vglue 0pt plus 1fill

\noindent
Prohlašuji, že jsem tuto bakalářskou práci vypracoval samostatně a výhradně
s~použitím citovaných pramenů, literatury a dalších odborných zdrojů.

\medskip\noindent
Beru na~vědomí, že se na moji práci vztahují práva a povinnosti vyplývající
ze zákona č. 121/2000 Sb., autorského zákona v~platném znění, zejména skutečnost,
že Univerzita Karlova v Praze má právo na~uzavření licenční smlouvy o~užití této
práce jako školního díla podle §60 odst. 1 autorského zákona.

\vspace{10mm}

\hbox{\hbox to 0.5\hsize{%
V ........ dne ............
\hss}\hbox to 0.5\hsize{%
Podpis autora
\hss}}

\vspace{20mm}


\newpage
%%% Povinná informační strana bakalářské práce

\vbox to 0.5\vsize{
\setlength\parindent{0mm}
\setlength\parskip{5mm}

Název práce:
Praktické datové struktury

Autor:
Michael Pokorný

Katedra:  % Případně Ústav:
Katedra aplikované matematiky % dle Organizační struktury MFF UK

Vedoucí bakalářské práce:
Mgr. Martin Mareš, PhD., Katedra aplikované matematiky

Abstrakt:
TODO
% TODO: v rozsahu 80-200 slov; nejedná se však o opis zadání bakalářské práce

Klíčová slova:
TODO % TODO: 3 až 5

\vss}\nobreak\vbox to 0.49\vsize{
\setlength\parindent{0mm}
\setlength\parskip{5mm}

Title:
Practical data structures

Author:
Michael Pokorný

Department:
Department of Applied Mathematics

Supervisor:
Mgr. Martin Mareš, PhD., Department of Applied Mathematics

Abstract:
TODO % v rozsahu 80-200 slov v angličtině; nejedná se však o překlad zadání bakalářské práce

Keywords:
TODO % 3 až 5 v angličtině

\vss}


\newpage

%%% Strana s automaticky generovaným obsahem diplomové práce. U matematických
%%% prací je přípustné, aby seznam tabulek a zkratek, existují-li, byl umístěn
%%% na začátku práce, místo na jejím konci.

\openright
\pagestyle{plain}
\setcounter{page}{1}
\tableofcontents

%%% Jednotlivé kapitoly práce jsou pro přehlednost uloženy v samostatných souborech
\chapter*{Úvod}
\addcontentsline{toc}{chapter}{Úvod}


\chapter{Název první kapitoly}

\section{Název první podkapitoly v první kapitole}

\section{Název druhé podkapitoly v první kapitole}

\subsection{Ukázka}
\label{ssec:ukazka}

\begin{itemize}
\item Logo Matfyzu vidíme na obrázku~\ref{fig:mff}.
\item Tato subsekce má číslo~\ref{ssec:ukazka}.
\end{itemize}

červeno-černý (krátká),
strana 16--22 (střední),
$45-44$ (minus),
a~toto je --- jak se asi dalo čekat --- vložená věta ohraničená dlouhými pomlčkami.
(Všimněte si, že jsme za \verb|a| napsali vlnovku místo mezery: to aby se
tam nemohl rozdělit řádek.)

\newtheorem{theorem}{Věta}
\newtheorem*{define}{Definice}	% Definice nečíslujeme, proto "*"

\begin{define}
{\sl Strom} je souvislý graf bez kružnic.
\end{define}

\begin{theorem}
Tato věta neplatí.
\end{theorem}

\begin{proof}
Neplatné věty nemají důkaz.
\end{proof}

\begin{figure}
	\centering
	\includegraphics[width=30mm]{img/logo.eps}
	\caption{Logo MFF UK}
	\label{fig:mff}
\end{figure}

\begin{itemize}
\item Model externí paměti

	CPU + malá cache + velká externí paměť.
	Zanedbáme všechny cache kromě poslední.

	Block size $B$ (v ext. paměti i v cachi),
	velikost cache $M$, tedy $M/B$ cache lines.

	Registry, L1, L2, RAM, disk, síť, ...
	Typicky L1 cache line = L2 cache line = 64B.
	\cite{Vit}

	\begin{itemize}
	\item Základní datové struktury: queue, stack
	\end{itemize}

\item Co je to slovník, operace
	\begin{itemize}
	\item Příklady použití
	\item Rozšíření: successor, predecessor
		\begin{itemize}
		\item Příklady použití
		\end{itemize}
	\item Rozšíření: multiset
		\begin{itemize}
		\item Příklady použití
		\end{itemize}
	\end{itemize}

\item Triviální implementace slovníku
	\begin{itemize}
	\item Triviální pole (unrolled linked list?)
	\item Triviální přímá adresace pro malé klíče

\item Hashování
	\begin{itemize}
		\item Seznam prvků v kolizi
		\item Linear probing (open addressing)
		\item Double hashing
		\item Cuckoo hashing
	\end{itemize}

\item Vyhledávací stromy
	\begin{itemize}
		\item B-stromy (a 2-3-4 stromy)
		\item Červenočerné stromy
		\item Splay stromy (TODO: zkusit vymyslet N-arni splay stromy)
		\item Van Emde-Boas uspořádání
		\item T-stromy (?)
	\end{itemize}

\item Speciality pro vnější paměť
	\begin{itemize}
		\item Directory hashing (extensible hashing, linear hashing)
		\item Directoryless hashing methods
		\item Tree EM structures
	\end{itemize}


\item Popis prostředí, kde měřím
\end{itemize}

\chapter*{Závěr}
\addcontentsline{toc}{chapter}{Závěr}


%%% Seznam použité literatury
%%% Seznam použité literatury je zpracován podle platných standardů. Povinnou citační
%%% normou pro diplomovou práci je ISO 690. Jména časopisů lze uvádět zkráceně, ale jen
%%% v kodifikované podobě. Všechny použité zdroje a prameny musí být řádně citovány.

\def\bibname{Seznam použité literatury}
\bibliographystyle{plain}
\bibliography{bibliography}

% \begin{thebibliography}{99}
% \addcontentsline{toc}{chapter}{\bibname}
% 
% %\bibitem{lamport94}
% %  {\sc Lamport,} Leslie.
% %  \emph{\LaTeX: A Document Preparation System}.
% %  2. vydání.
% %  Massachusetts: Addison Wesley, 1994.
% %  ISBN 0-201-52983-1.
% 
% \end{thebibliography}


%%% Tabulky v diplomové práci, existují-li.
\chapwithtoc{Seznam tabulek}

%%% Použité zkratky v diplomové práci, existují-li, včetně jejich vysvětlení.
\chapwithtoc{Seznam použitých zkratek}

%%% Přílohy k diplomové práci, existují-li (různé dodatky jako výpisy programů,
%%% diagramy apod.). Každá příloha musí být alespoň jednou odkazována z vlastního
%%% textu práce. Přílohy se číslují.
\chapwithtoc{Přílohy}

\openright
\end{document}
